\documentclass[10pt]{article}
\usepackage[margin=20mm]{geometry}
\usepackage{amsmath,amssymb,mathtools,fontspec,unicode-math,xurl,hyperref}
\IfFontExistsTF{DejaVu Serif}{\setmainfont{DejaVu Serif}}{\setmainfont{Latin Modern Roman}}
\setmathfont{latinmodern-math.otf}
\setlength{\parindent}{0pt}\setlength{\parskip}{6pt}\sloppy
\title{Scientific initial review M25-0008
The Alperin–McKay conjecture for the prime 2}
\author{}
\date{26 October 2026}
\begin{document}
\maketitle
=3em Object and disposition. This review concerns Lucas Ruhstorfer’s arXiv:2204.06373v1, submitted 13 April 2022, with 54 PDF pages and the earlier title printed above. It identifies two Major issues in the written proof chain and five Minor groups of local errors. No counterexample to the Alperin–McKay conjecture or Brauer’s height-zero conjecture is asserted. Proposed repairs are distinguished from adopted conclusions.\par
The frozen journal identity is The Alperin–McKay and Brauer’s height zero conjecture for the prime 2, Annals of Mathematics 201 (2025), 379–457, DOI 10.4007/annals.2025.201.2.2. The supplied identity recovery links this DOI and arXiv identifier in the same zbMATH record 8014492 and supplies corroborating references. This establishes correspondence, not equality of texts. The 2025 Version of Record was not obtained; no finding below is attributed to that uninspected version. The earlier manuscript already states a height-zero corollary on page 1.\par
\section*{Scope, version, and method}
All 54 source PDF pages were read. The review followed definitions, quantifiers, block and character parametrizations, Sylow and root calculations, extension maps, local-to-global criteria, and the exceptional-case assembly. Twelve source pages were additionally inspected as rendered images: 7, 9, 17, 19, 31, 34, 36, 44, 46, 47, 51, and 52. The 13 manifest file bindings were checked by size and SHA256. No author TeX source is supplied for this item; static inspection is limited to the actual PDF and the reviewer’s own sources. No author code was run.\par
The arXiv history checked for this review lists v1 only. The publisher page records receipt on 14 April 2022, revision on 7 August 2024, acceptance on 13 March 2024, and online publication on 12 March 2025; the unusual displayed ordering of revision and acceptance is retained as metadata, not silently corrected. Bounded searches did not identify an article-specific erratum; this is not proof that none exists. Remark 4.6, page 17, corrects an assertion in Cabanes’s 2018 work, and Lemma 7.4, page 26, supplies an omitted centralizer in Malle’s 2018 work. Those are historical issues acknowledged in the reviewed manuscript, not new defects assigned to it.\par
Primary dependency checks included Ruhstorfer’s local criterion [Ruh], Brough–Späth’s extension construction [BS], and the original character tables of Simpson–Frame [SF]. Classification, general Deligne–Lusztig theory, all tables in Kessar–Malle, and every cited equivalence theorem have not been independently re-proved. These are explicit limits on coverage, not assertions that such inputs are false.\par
\section*{Major issue F01: the local subgroup hypothesis cannot hold}
Location. Theorem 9.3, page 31, first sentence; its application in Theorem 11.3, pages 41–42; compare Lemma 6.6(a), page 24, and Definition 9.1, page 31.\par
Claim and complete check. Theorem 9.3 demands a proper \(\operatorname{Aut}(G)_b\)-stable subgroup \(M\) containing \(N_G(P)\). Every inner automorphism fixes a block \(b\): a block idempotent is central in the group algebra. Thus \[\operatorname{Inn}(G)\leq\operatorname{Aut}(G)_b,\qquad
 M\text{ is }\operatorname{Aut}(G)_b\text{-stable}\ \Longrightarrow\ M\lhd G.\] In the quasisimple groups at issue every proper normal subgroup is central. Since \(P\leq N_G(P)\leq M\) and the relevant defect group is noncentral, the demanded proper subgroup does not exist. This is an incompatibility of the stated hypotheses, not a counterexample to an implication with impossible hypotheses.\par
Lemma 6.6(a) proves the different assertion that \(N_G(L)\) is stable under \(\operatorname{Aut}(G)_{P,b}\). Theorem 11.3 explicitly selects \(M=N_G(L)\) and says that Lemma 6.6 verifies the assumptions of Theorem 9.3. It does not verify the printed \(\operatorname{Aut}(G)_b\)-stability condition. The same overly large stabilizer reappears in the argument of Corollary 14.3 on page 51.\par
Counterevidence and repair. Definition 9.1 uses the appropriate group \(N_{\operatorname{Aut}(H)}(D,b)\), which preserves the chosen defect group. Moreover, Theorem 2.1 of the primary source [Ruh], page 8 of arXiv:2010.04499v2, is formulated with \(M=N_G(Q)\), for \(Q\) characteristic in a defect group, rather than with the impossible invariance under all inner automorphisms. These strongly support a missing stabilizer subscript. The repair must state the correct local criterion, verify its normalizer/characteristic-subgroup formulation or the required generalization to \(N_G(L)\), and then check its remaining character and block conditions. Merely retaining the printed vacuous criterion does not do this.\par
Core impact and judgment. The affected application is the final iAM verification in Theorem 11.3, used in the noncyclic-stabilizer branch of Theorem 14.6; Lemma 9.6 and other subsequent uses also require the corrected criterion. This is a demonstrated invalid derivation on the main proof path, hence Major under the supplied standard, even though a local notational repair is plausible. The main conjectures are not refuted.\par
\section*{Major issue F02: a factor-swap cannot be in the extension kernel}
Location. Lemma 12.6, last paragraph of page 44 and first paragraph of page 45; Remark 12.4, page 44; application in Theorem 12.8, page 45.\par
Here \(L_0=\operatorname{SL}_3(q)\times\operatorname{SL}_3(q)\), \(W_{(d)}\simeq S_3\), its subgroup of order three centralizes \(L_0\), and its involutions interchange the two simple factors. The proof claims an extension of \(\psi_0\) to \(L_0\rtimes W_{(d)}\) with all of \(W_{(d)}\) in its kernel, and uses it to construct the character denoted \(\psi_0\Gamma(\nu)\).\par
Let \(K=\operatorname{SL}_3(q)\) and let \(\tau\) swap the factors. Any representation \(D\) with \(D(\tau)=I\) must satisfy \[D(a,1)=D\bigl(\tau(a,1)\tau^{-1}\bigr)=D(1,a)\quad(a\in K).\] The two commuting copies of \(K\) therefore have the same image, and that image is abelian. Since \(K\) is perfect for the cases here, \(D|_{K\times K}\) is trivial. Thus no nontrivial \(\psi_0\) of the required kind can have such an extension.\par
Equivalently, for \(\psi_0=\chi\otimes\chi\), the swap must act by a tensor permutation: \[F(v\otimes w)=w\otimes v,\qquad
 F(\rho(a)\otimes I)F^{-1}=I\otimes\rho(a).\] Replacing \(F\) by \(I\) forces \(\rho(a)\) to be scalar for every \(a\). The obstruction is not restricted to a toy representation. The relevant \(A_2\) series is nontrivial on the simple factors; for example at the allowed split parameter \(q=13\), its exceptional constituent has degree \((q-1)^2(q+1)/3=672\), as also described in Section 12.4. Its tensor square cannot satisfy the claimed kernel condition. An explicit \(4\)-by-\(4\) matrix calculation in the independent script checks the swap relation and the failure of the identity operator, while the argument above proves the general obstruction.\par
Counterevidence and repair. The canonical tensor-permutation extension does exist. Its kernel contains \(C_{W_{(d)}}(L_0)\), the subgroup of order three, not the whole permutation group. Proposition 4.2 of [BS], page 16, requires an actual extension representation \(D'\); its formula (4.1) retains the factor \(D'(\rho(v))\). It does not permit dropping a nontrivial swap operator. Replace the false kernel assertion, carry the tensor-permutation action through the construction on \(L_0V_{(d)}\), and check the \(E'\)-stabilizer and induction steps on page 45 with that action retained.\par
Core impact and judgment. The displayed construction supplies the equivariant extension map used in Theorem 12.8 for the \(A_2(\varepsilon q^3).3\) cases, part of the final exceptional-group coverage. Its printed justification cannot be executed. This is Major; the existence of a plausible corrected extension is counterevidence against claiming that Lemma 12.6 itself, or the main result, is false.\par
\section*{Minor issues with explicit checks}
\subsection*{F03: unequal orbit lengths in Theorem 1.4(a)}
Location. Page 7, Theorem 1.4(a) and the last paragraph of its proof.\par
Take \(\mathbf G=\operatorname{SL}_2\) over \(\mathbb F_5\), with dual \(\operatorname{PGL}_2\), and let \(s\) be the image of \(\operatorname{diag}(1,-1)\). Here \(Z_F=C_2\), \(C^\circ_{\mathbf G^*}(s)\) is a split torus, and \(A(s)=C_2\). The torus has only the trivial unipotent character, so its \(A(s)\)-orbit has size one.\par
For the regular embedding in \(\operatorname{GL}_2(5)\), let \(\eta\) be the nontrivial quadratic character of \(\mathbb F_5^\times\). The principal series induced from \((1,\eta)\) is irreducible of degree six: the two Bruhat double cosets give endomorphism dimension one when the two inducing characters differ. Twisting by \(\eta\circ\det\) interchanges the two inducing characters and fixes this representation. These are exactly the two determinant twists fixing it. Cyclic Clifford theory therefore gives two distinct degree-three constituents on \(\operatorname{SL}_2(5)\), interchanged by diagonal conjugation. Their \(Z_F\)-orbit has size two. Thus the corresponding orbit lengths are \(2\) and \(1\), not equal.\par
For \(a=|A(s)|\) and \(r\) the unipotent orbit length, the relevant cyclic Clifford relations are \[\#(\text{diagonal orbit})=a/r,\qquad
 \chi(1)=\widetilde\chi(1)\,r/a,\qquad
 \psi(1)=r\,\widetilde\psi(1).\] The printed formulas divide \(\widetilde\chi(1)\) by \(r\) and multiply \(\widetilde\psi(1)\) by \(a/r\) instead. In this example they would give \(6\) and \(2\), whereas the actual degrees are \(3\) and \(1\). The script enumerates \(\operatorname{GL}_2(5)\) (480 elements), \(\operatorname{SL}_2(5)\) (120), and the centralizers (orders 16 upstairs and 8 projectively), corroborating the degree indices.\par
Counterevidence/impact. The quotient of the two corrected degree formulas still yields the degree formula in part (c). Lemma 2.3 constructs its later correspondence through additional Clifford arguments. No use of the false equality of orbit lengths that defeats the final AM count was established. Correct the orbit statement and its proof; classify this verified auxiliary error as Minor, without promoting it to a counterexample to part (c).\par
\subsection*{F04: the Cabanes tests have overstrong quantifiers}
Location. Proposition 4.2 and Lemma 4.4, page 17; compare the quotient-involution argument in Theorem 4.10, pages 18–20.\par
Proposition 4.2 claims necessity as well as sufficiency of the absence of nontrivial quadratic cosets. For necessity take \(H=S_4\) and \(A=V_4\). This \(A\) is the unique maximal abelian normal subgroup. Indeed, for any abelian normal \(B\), minimal normality of \(A\) gives \(B\cap A=1\) or \(A\). In the second case \(B\leq C_H(A)=A\); in the first, \([B,A]\leq B\cap A=1\) gives the same conclusion. Yet a transposition induces a nontrivial involution on \(A\simeq\mathbb F_2^2\), and in characteristic two \((h-I)^2=h^2-2h+I=0\). This is the forbidden quadratic coset in a group satisfying the printed conclusion. The independent enumeration checks the abelian normal closures and the action.\par
The sufficient direction survives: if an abelian normal \(B\) is not contained in \(A\), choose \(h\in B\setminus A\). Then \([h,A]\subseteq B\cap A\) centralizes \(h\), making the coset quadratic.\par
Lemma 4.4 incorrectly asks for an actual involution in \(H\setminus A\). Take \(H=Q_8\) and \(A=\langle i\rangle\). The quotient is \(C_2\), there is no involution outside \(A\), and \(H\) has three maximal abelian normal subgroups of order four. Its nontrivial coset is involutory and acts quadratically on \(A\), but its representatives have order four. The correct obstruction is an involutory coset acting on \(A\). For such a coset the action squares to one, and the quadratic condition becomes \([h,t]^2=1\) for all \(t\in A\). The further implication from equality of squares to \({}^ht\in\{t,t^{-1}\}\) is also invalid for general abelian \(A\) and is unnecessary.\par
Counterevidence/impact. The proof of Theorem 4.10 tests involutions in the Weyl quotient, which is the corrected type of test; it does not depend on finding order-two lifts in the Sylow group. Thus these two finite counterexamples do not refute its Sylow conclusion. Retain sufficiency, use cosets, and remove the unnecessary equality-of-squares step. This is a Minor group of false auxiliary assertions, with no demonstrated failure of the core Sylow conclusion.\par
\subsection*{F05: the three \(A_2\) subgroups do not share pairwise a \(C_3\)}
Location. Section 10.3, page 34, last paragraph; simply connected torus convention on page 33.\par
The root subgroups \(G_1,G_2,G_3\) have bases \(\{\alpha_1,\alpha_3\}\), \(\{\alpha_5,\alpha_6\}\), and \(\{\alpha_0,\alpha_2\}\). The text asserts \(G_i\cap G_j=Z\) of order three for every \(i\ne j\). But \(G_1\) and \(G_2\) commute and have independent coordinate subtori in the simply connected torus, whose simple coroots form a basis. Their intersection is central in both, hence lies in those subtori, and is trivial. The same calculation applies to the other pairs.\par
There is a genuine order-three kernel in the product map. With the order of simple coroots used in the paper, its lattice matrix has columns \[e_1,\ e_3,\ e_5,\ e_6,\ -(1,2,2,3,2,1)^t,\ e_2\] and determinant of absolute value three. The kernel involves all three factors diagonally; it is not a common pairwise intersection. Exact determinant arithmetic and enumeration on the three-torsion subtori give pairwise intersection sizes \(1,1,1\) and product-kernel size \(3\).\par
Counterevidence/impact. The correct central-product description and rank are close to the intended statement. The discrepancy concerns odd central torsion; the subsequent explicit two-local calculations do not thereby become false. No necessary later use of a nontrivial pairwise intersection was identified. Replace the assertion by the product-kernel formulation. Minor.\par
\subsection*{F06: central characters must be normalized by degree}
Location. Lemma 12.9(a),(b), pages 46–47.\par
Part (a) prints \(\operatorname{Res}_{Z(U)}^U(\phi)\phi(1)=\operatorname{Res}_{Z(U)}^H(\chi)\chi(1)\). At the identity this would require \(\phi(1)^2=\chi(1)^2\), even when restriction has nontrivial multiplicity. All hypotheses are met by \[H=S_3\times C_3,\quad H_0=U=C_3,\quad Z=1,\quad
 \chi=\chi_{\rm std}\otimes1,\quad\phi=1_U.\] Here \(\chi(1)=2\), \(\phi(1)=1\), and the multiplicity is two. The printed identity is \(1=4\).\par
Schur’s lemma gives the correct statement \[\frac{\operatorname{Res}_{Z(U)}^U\phi}{\phi(1)}
   =\frac{\operatorname{Res}_{Z(U)}^H\chi}{\chi(1)}.\] Part (b) must use the corresponding normalized central character on \(Z(U)Z\). If \(\lambda_\chi\) denotes the scalar central character of \(\chi\), the extension is explicitly \(\widehat\rho(uz)=\rho_\phi(u)\lambda_\chi(z)\); agreement on \(U\cap Z\) proves it well defined, and determines it uniquely.\par
Counterevidence/impact. The proof in words already establishes equality of scalar central characters, and Proposition 12.10 needs that normalized assertion. The correction is fully supplied here; its odd-multiplicity argument is checked separately below. Thus Minor, not evidence that the claimed odd multiplicity is false.\par
\subsection*{F07: the defining union for \(\mathcal E_\ell(G,s)\) omits \(s\)}
Location. Section 1.6, page 9, first display.\par
The union is printed with \(\mathcal E(G,t)\) instead of \(\mathcal E(G,st)\). Take \(\mathbf G=\mathbb G_m\) over \(\mathbb F_7\), so \(G=C_6\), \(\ell=2\), and \(s\) has exponent \(2\) and order \(3\) in its dual. The displayed union selects characters of exponents \(\{0,3\}\), whereas the series associated with \(s\) is \(\{2,5\}\). These are different \(2\)-blocks. In a torus the printed expression loses its dependence on \(s\) altogether.\par
Counterevidence/impact. Lemma 2.3, page 11, explicitly uses the products \(st\). Thus this is a local missing factor, not an independently demonstrated failure of the later parametrization. Insert \(s\) in the defining union. Minor.\par
\section*{Positive checks and resolved concerns}
\subsection*{An independent verification of the character-table calculation}
Proposition 12.10, page 47, invokes CHEVIE for an odd-multiplicity claim. The following derivation checks its split \(\operatorname{SL}_3(q)\) calculation directly from the original tables [SF], Tables 1a–1b, with \(q\equiv1\pmod3\). It uses no author code and no CHEVIE execution.\par
Write \(r=q-1\), \(s=q+1\), \(t=q^2+q+1\), \(r'=r/3\), and \(D=r^2s/3\). The six exceptional characters of degree \(D\) have indices \(\chi_{u,v}\), \(u=0,1,2\), \(v=1,2\). With \(\omega^3=1\), their relevant nonzero values are: \(D\omega^{vk}\) on the three central classes, \(-r'\omega^{vk}\) on the three transvection classes, \((q\delta_{\ell u}-r')\omega^{vk}\) on the nine regular-unipotent classes, and \(\omega^{vk}\) on irreducible-torus classes. The corresponding centralizer orders are \(q^3r^2st\), \(q^3r\), \(3q^2\), and \(t\); the last class family has \((t-3)/3\) members.\par
For \(\chi_{w,2v}\) (with the second index reduced modulo three), central phases cancel in the scalar product. Hence \[\begin{aligned}
 \langle\chi_{u,v}^2,\chi_{w,2v}\rangle
 &=C(q)+q^{-2}\sum_{\ell=0}^2
 (q\delta_{\ell u}-r')^2(q\delta_{\ell w}-r'),\\
 C(q)&=\frac{3D^3}{q^3r^2st}
       -\frac{3(r')^3}{q^3r}
       +\frac{t-3}{3t}
       =\frac{q^3-1}{9q^2}.\end{aligned}\] The remaining sum is \((2q^3+6q^2+1)/9\) if \(w=u\), and \((-q^3+1)/9\) otherwise. Therefore \[\boxed{\ \langle\chi_{u,v}^2,\chi_{w,2v}\rangle
       =\begin{cases}(q+2)/3,&w=u,\\0,&w\ne u.\end{cases}\ }\] Characters with the incompatible central character have scalar product zero. For odd \(q\equiv1\pmod3\), \((q+2)/3\) is odd. This proves the required uniqueness and parity in this family for all such \(q\), rather than extrapolating from examples. The script verifies the rational polynomial identity coefficient by coefficient and evaluates all nine scalar products for twelve values of \(q\), including \(q=13\) with degree \(672\) and multiplicity \(5\). One even \(q\) is included only as an algebraic diagnostic, not as an odd-characteristic case.\par
For the subfield case on the same page, a regular element of order \(q^2+q+1\) remains in the irreducible-torus class family over \(\mathbb F_{q^2}\). The embedding exponent \(q^2-q+1\) is \(1\) modulo three. The table therefore gives a third root of unity for the relevant exceptional character value. Among constituents nonzero at that element and with its nontrivial central character, the exceptional and generic cuspidal degrees have the required common two-part; the other relevant characters vanish. If all their multiplicities were even, a third root of unity divided by two would be an algebraic integer, which is impossible (its norm in \(\mathbb Q(\omega)\) is \(1/4\)). This verifies the parity mechanism used there. No blanket verification of every unitary table entry or every Ennola substitution is claimed.\par
\subsection*{Root, group-action, and case-assembly checks}
The independent program constructs all 72 \(E_6\) roots, identifies the 36 positive roots, and verifies the coefficient assertion of Lemma 4.8 for that root system. Orthogonality to the tori in Section 12.2 gives 12 and 14 roots, agreeing with the claimed \(A_2^2\) and \(A_2^2A_1\) Levi types. These checks also support the lattice calculation in F05.\par
The statement about an odd group acting on a power-of-two set in the argument of Lemma 2.3 needs its actual context: here the relevant odd component group is at most \(C_3\). Since \(3\) does not divide a power of two, a fixed point exists; the commuting transitive diagonal action then propagates fixedness. An arbitrary odd group acting on an arbitrary power-of-two set would not justify the same shortcut. No separate defect is assigned after restoring the context.\par
Similarly, the coprime-orbit argument in Proposition 8.3 uses the commuting action supplied by the cyclic outer automorphism setup. A permutation commuting with a transitive abelian action acts as a translation on the orbit; coprimality then forces the required trivial action. Coprimality alone, without this structure, would not suffice.\par
The final cyclic-stabilizer sentence in Theorem 14.6 cites Proposition 14.4, whose statement excludes two centralizer types. The already proved Theorem 12.8 and Lemma 14.5 address those excluded cases. The final assembly should cite them explicitly. This is not recorded as a new missing-case obstruction.\par
\section*{Judgment, repairs, and adoption}
The review establishes failures of two printed core derivations and five groups of local assertions. It does not establish that the main conclusions are false. A repair proposed in this report is not silently substituted into the reviewed manuscript.\par
Item | Judgment | Repair or positive evidence | Adoption / remaining work\par
F01 | Major | Use defect-preserving local criterion | Verify corrected criterion and each application\par
F02 | Major | Retain tensor-permutation operator | Recheck extension and equivariance construction\par
F03 | Minor | Reciprocal orbit relation | Correct statement and degree calculation\par
F04 | Minor | Sufficient test; involutory cosets | Revise auxiliary statements and lifts\par
F05 | Minor | Product kernel of order three | Replace pairwise-intersection assertion\par
F06 | Minor | Divide central restrictions by degrees | Local correction proved here\par
F07 | Minor | Use the product \(st\) | Local correction consistent with later use\par
Positive | Adopt locally | Exact odd-multiplicity formula; root checks | Does not certify all external inputs\par
Version | Limited | Identity correspondence established | 2025 journal text remains uninspected\par
\section*{Supporting checks}
The input PDF SHA256 is\par
> 00000TOKEN\par
The manifest SHA256 is\par
> 00001TOKEN\par
The identity-recovery receipt SHA256 is\par
> 00002TOKEN\par
> 00004TOKEN\par
was executed by this reviewing assistant with Python, returned exit code zero, and wrote `independent\_checks.json`. It uses standard-library exact arithmetic and finite enumerations. Its output records the \(S_4\), \(Q_8\), \(\operatorname{GL}_2(5)\), lattice, root, tensor, normalized-character, torus-series, and character-table checks. These are targeted independent checks, not exhaustive verification of a 54-page proof. No failed scientific computation is hidden.\par
The downloaded primary Simpson–Frame PDF, bounded to a small dependency file, has SHA256\par
> 00006TOKEN\par
Its relevant table pages were actually inspected as images. Retrieval status, URL, byte count, and hash are recorded in `dependencies/SF1973\_receipt.json`; the derivation above is the reviewer’s own.\par
The exact report source is `reports/0008.tex`; its actual compiled output is `reports/Pdf/0008.pdf`. The PDF’s hash is deliberately external to this PDF to avoid a self-referential checksum. The render log and page-by-page inspection record bind the checked images to that PDF.\par
The authorized Tectonic 0.17.0 executable was independently checked and invoked with `--untrusted`; shell escape and author-code execution were not enabled. Task-local cache and temporary use were bounded, and both Windows disks were checked against the required free-space floor.\par
\section*{Supporting checks}
All scientific reading, proof checks, independent computations, static inspection, compilation, and actual page inspection for this report were performed by the visible reviewing assistant. No scientific work was delegated. The dispatcher’s observed account and interface were Minhao Liu’s Group — Academic Researchers, Chat mode, webpage 6Pro. This is attributed to the dispatcher, not claimed as a new independent UI observation. No model switch was requested or performed.\par
> 00018TOKEN\par
9 L. Ruhstorfer, The Alperin–McKay conjecture for the prime 2, arXiv:2204.06373v1 (2022), the exact reviewed manuscript. https://arxiv.org/abs/2204.06373v1 Annals of Mathematics, publisher metadata for DOI 10.4007/annals.2025.201.2.2. Identity and revision metadata only; journal PDF uninspected. https://annals.math.princeton.edu/2025/201-2/p02 L. Ruhstorfer, Jordan decomposition for the Alperin–McKay conjecture, primary version arXiv:2010.04499v2, especially Theorem 2.1. https://arxiv.org/abs/2010.04499v2 J. Brough and B. Späth, On the Alperin–McKay conjecture for simple groups of type A, arXiv:1901.04591v1, Proposition 4.2 and formula (4.1). https://arxiv.org/abs/1901.04591v1 W. A. Simpson and J. S. Frame, The character tables for \(\operatorname{SL}(3,q)\), \(\operatorname{SU}(3,q^2)\), \(\operatorname{PSL}(3,q)\), \(\operatorname{PSU}(3,q^2)\), Canad. J. Math. 25 (1973), 486–494, Tables 1a–1b. https://doi.org/10.4153/CJM-1973-049-7\par
\end{document}
